\section{Conclusion}
\label{sec:conclusion}

This review reframes VLA attacks as a closed-loop robotics problem. The field now has direct evidence for adversarial patches, black-box instruction attacks, action freezing, action-level and long-horizon backdoors, transferable and 3D object attacks, reasoning-chain hijacking, trajectory redirection, and membership inference. That evidence is important, but it is not yet uniformly authoritative: it is recent, often preprint-heavy, concentrated around OpenVLA and LIBERO, and uneven in physical recovery, HIL, controller reporting, and adaptive evaluation.

The main conclusion is therefore evidence-aware but mechanism-driven. Direct VLA studies support claims about demonstrated VLA attack surfaces. Embodied-adjacent studies motivate hypotheses about planners, tools, sensors, middleware, and HRI that still need VLA validation. Contextual studies explain mechanisms and evaluation design, not VLA vulnerability by themselves. A shared typed definition of attack success is necessary for this field: task failure, targeted action, safety violation, authorization failure, privacy leakage, persistence, and recovery failure are different predicates and should not be collapsed into one ASR. Progress now depends on cross-representation transfer studies, physical and HIL protocols, privacy-aware evaluation, adaptive attacks against deployed guardrails, and responsible artifact release suited to embodied systems.
